\documentclass[11pt,letterpaper]{article}
\usepackage{math-template}

% Edit these fields. Empty optional fields are hidden.
\newcommand{\StudentName}{Your Name}
\newcommand{\CourseName}{MATH 101 --- Introduction to Analysis}
\newcommand{\AssignmentTitle}{Assignment 1}
\newcommand{\DueDate}{October 15, 2026}
\newcommand{\StudentNumber}{}
\newcommand{\TutorialSection}{}

\begin{document}
\makeassignmentheader

\begin{problem}{Triangle inequality}{triangle}
Show that $\abs{x+y}\leq\abs{x}+\abs{y}$ for all $x,y\in\R$.
\end{problem}
\begin{solution}
Adding $-\abs{x}\leq x\leq\abs{x}$ and
$-\abs{y}\leq y\leq\abs{y}$ gives
\[
  -(\abs{x}+\abs{y})\leq x+y\leq\abs{x}+\abs{y}.
\]
The claim follows from the definition of absolute value.
\end{solution}

% Each box inside this group is numbered 2(a), 2(b), ...
\begin{subproblems}
\begin{problem}{A convergent sequence}{limit}
Prove that $a_n=1/n$ converges to $0$.
\end{problem}
\begin{solution}
Let $\eps>0$. Choose $N\in\N$ with $N>1/\eps$.
For every $n\geq N$, we have $\abs{a_n-0}=1/n\leq1/N<\eps$.
Hence $a_n\to0$.
\end{solution}

\begin{problem}{A shifted sequence}{shift}
Deduce that $b_n=3+1/n$ converges to $3$.
\end{problem}
\begin{solution}
Since $\abs{b_n-3}=1/n$, the same estimate as in
Problem~\ref{p:limit} proves the claim.
\end{solution}
\end{subproblems}

\begin{problem}{An integral}{integral}
Evaluate $\int_0^1 x^2\dd x$.
\end{problem}
\begin{solution}
By the fundamental theorem of calculus,
\[
  \int_0^1 x^2\dd x=\left[\frac{x^3}{3}\right]_0^1=\frac13.\qedhere
\]
\end{solution}
\end{document}
